\documentclass{article}

% ============================================================
% MA 213 In-Class Worksheet Template
% Student-facing worksheet for lecture activities.
% ============================================================

\usepackage[margin=1in]{geometry}
\usepackage{amsmath}
\usepackage{xcolor}
\usepackage{parskip}
\usepackage{array}
\usepackage{enumitem}
\usepackage{booktabs}

\newcommand{\coursenumber}{MA 213}
\newcommand{\weekplaceholder}{14}
\newcommand{\lectureplaceholder}{33}
\newcommand{\lecturetitle}{Spam or Not Spam? Think/Pair/Share}

\newcommand{\nameblank}{\rule{1.75in}{0.4pt}}
\newcommand{\idblank}{\rule{1.55in}{0.4pt}}
\newcommand{\shortblank}{\rule{1in}{0.4pt}}
\newcommand{\longblank}{\rule{0.93\linewidth}{0.4pt}}

\begin{document}

\begin{center}
  {\LARGE \coursenumber{} Week \weekplaceholder{}: Lecture \lectureplaceholder{}}\\
  {\large \lecturetitle{}}
\end{center}

\begin{enumerate}[label=\arabic*.,leftmargin=*,itemsep=.7em]
  \item \textbf{Your name:} \nameblank \hfill \textbf{BU ID:} \idblank
\end{enumerate}

\textbf{Big idea:} You just watched a two-hypothesis posterior get computed for a disease test, and then updated \emph{again} using a second positive test as new evidence. Now you'll run that same two-step process yourself, on a different two-hypothesis problem: is an email spam?

\section*{The Setup}

An email either is \textbf{Spam} or \textbf{Not Spam}. Based on your inbox history, your prior belief is
\[
P(\text{Spam}) = 0.20, \qquad P(\text{Not Spam}) = 0.80.
\]

\section*{Step 1: Think (5 mins)}

\textbf{(a)} The email contains the word ``free.'' Historically,
\[
P(\text{``free''} \mid \text{Spam}) = 0.30, \qquad P(\text{``free''} \mid \text{Not Spam}) = 0.05.
\]
Use Bayes' rule to compute $P(\text{Spam} \mid \text{``free''})$.
\vspace{0.4in}

\textbf{(b)} The same email \emph{also} contains the word ``winner.'' Historically (and independently of ``free,'' given the class),
\[
P(\text{``winner''} \mid \text{Spam}) = 0.25, \qquad P(\text{``winner''} \mid \text{Not Spam}) = 0.02.
\]
\textbf{Yesterday's posterior becomes today's prior:} use your answer from (a) as the new prior, and update again on ``winner.'' What is $P(\text{Spam} \mid \text{``free,'' ``winner''})$?
\vspace{0.4in}

\section*{Step 2: Pair (4 mins)}

\begin{enumerate}[label=\arabic*.,leftmargin=*,itemsep=.7em, start=2]
  \item \textbf{Partner's name:} \nameblank \hfill \textbf{BU ID:} \idblank
\end{enumerate}

Compare your answers to (a) and (b) with your partner. Then discuss:

\begin{itemize}
  \item Would you have gotten the same final posterior in part (b) if you'd updated on ``winner'' first and ``free'' second? Why or why not?
  \item \textbf{Bonus:} suppose the email contained \emph{neither} ``free'' nor ``winner.'' Without recomputing, would you expect $P(\text{Spam} \mid \text{data})$ to be higher or lower than your prior of 0.20? Why?
\end{itemize}

\textbf{Our thoughts:}
\vfill

\end{document}
